% !TEX program = pdflatex
% Projeto de pesquisa para edital ou projeto de tese. Ajuste as seções ao
% que o edital pede e ao limite de páginas.
\documentclass[12pt,a4paper]{article}
\usepackage[T1]{fontenc}
\usepackage[utf8]{inputenc}
\usepackage[brazilian]{babel}
\usepackage{lmodern}
\usepackage{microtype}
\usepackage[margin=2.5cm]{geometry}
\usepackage{setspace}\onehalfspacing
\usepackage{amsmath}
\usepackage{booktabs,tabularx}
\usepackage{enumitem}
\usepackage[authoryear,round]{natbib}
\usepackage[hidelinks]{hyperref}

\title{{{TITLE}}}
\author{{{AUTHOR}}\\{{INSTITUTION}}}
\date{{{DATE}}}

\begin{document}
\maketitle

\begin{abstract}
\noindent Em até 250 palavras: o problema, por que importa, o que o projeto faz, com que dados e método, e o que entrega.

\medskip
\noindent\textbf{Palavras-chave:} palavra um; palavra dois; palavra três.
\end{abstract}

\section{Problema e justificativa}
Qual é a pergunta, por que ainda não tem resposta e quem ganha com ela. A literatura entra aqui só para mostrar a lacuna \citep{card1994minimum}.

\section{Objetivos}
\paragraph{Objetivo geral.} Uma frase.
\paragraph{Objetivos específicos.}
\begin{enumerate}[nosep]
  \item Primeiro objetivo, verificável.
  \item Segundo objetivo.
\end{enumerate}

\section{Dados e metodologia}
Fontes de dados, unidade de observação, período e estratégia de identificação.
\begin{equation}
  y_{it} = \beta\, D_{it} + \alpha_i + \gamma_t + \varepsilon_{it}
\end{equation}

\section{Resultados esperados e impacto}
Produtos (artigos, bases de dados, notas técnicas) e quem os usa.

\section{Cronograma}
\begin{table}[h]
  \centering
  \begin{tabular}{lcccc}
    \toprule
    Atividade & Sem.~1 & Sem.~2 & Sem.~3 & Sem.~4 \\
    \midrule
    Montagem dos dados & $\bullet$ & $\bullet$ & & \\
    Estimação & & $\bullet$ & $\bullet$ & \\
    Redação e submissão & & & $\bullet$ & $\bullet$ \\
    \bottomrule
  \end{tabular}
\end{table}

\section{Orçamento}
\begin{table}[h]
  \centering
  \begin{tabularx}{0.8\textwidth}{Xr}
    \toprule
    Item & Valor (R\$) \\
    \midrule
    Bolsa de iniciação científica (12 meses) & 0,00 \\
    Software e computação & 0,00 \\
    Participação em congresso & 0,00 \\
    \midrule
    \textbf{Total} & \textbf{0,00} \\
    \bottomrule
  \end{tabularx}
\end{table}

\section{Equipe}
Quem faz o quê.

\bibliographystyle{aer}
\bibliography{references}

\end{document}
